\documentclass[10pt,a4paper]{amsart}
\usepackage{amsmath,amssymb,amsthm,graphicx}
% Retained minimal preamble for the Stacks Project chapter more-morphisms.tex (More on Morphisms).
% Verbatim excerpt of the chapter preamble: Xy-pic setup (native preamble lines 14-18) and the native
% theorem declarations with the mathematical macros (native preamble lines 154-217). Copied mechanically;
% nothing added, removed or redefined. The wrapper class and the amsmath/amssymb/amsthm/graphicx base
% packages are supplied by the assembler, so the class/externaldocument part of the native preamble is omitted.
% For commutative diagrams we use Xy-pic
\usepackage[all]{xy}

% We use 2cell for 2-commutative diagrams.
\xyoption{2cell}
% Theorem environments.
%
\theoremstyle{plain}
\newtheorem{theorem}[subsection]{Theorem}
\newtheorem{proposition}[subsection]{Proposition}
\newtheorem{lemma}[subsection]{Lemma}

\theoremstyle{definition}
\newtheorem{definition}[subsection]{Definition}
\newtheorem{example}[subsection]{Example}
\newtheorem{exercise}[subsection]{Exercise}
\newtheorem{situation}[subsection]{Situation}

\theoremstyle{remark}
\newtheorem{remark}[subsection]{Remark}
\newtheorem{remarks}[subsection]{Remarks}

\numberwithin{equation}{subsection}

% Macros
%
\def\lim{\mathop{\mathrm{lim}}\nolimits}
\def\colim{\mathop{\mathrm{colim}}\nolimits}
\def\Spec{\mathop{\mathrm{Spec}}}
\def\Hom{\mathop{\mathrm{Hom}}\nolimits}
\def\Ext{\mathop{\mathrm{Ext}}\nolimits}
\def\SheafHom{\mathop{\mathcal{H}\!\mathit{om}}\nolimits}
\def\SheafExt{\mathop{\mathcal{E}\!\mathit{xt}}\nolimits}
\def\Sch{\mathit{Sch}}
\def\Mor{\mathop{\mathrm{Mor}}\nolimits}
\def\Ob{\mathop{\mathrm{Ob}}\nolimits}
\def\Sh{\mathop{\mathit{Sh}}\nolimits}
\def\NL{\mathop{N\!L}\nolimits}
\def\CH{\mathop{\mathrm{CH}}\nolimits}
\def\proetale{{pro\text{-}\acute{e}tale}}
\def\etale{{\acute{e}tale}}
\def\QCoh{\mathit{QCoh}}
\def\Ker{\mathop{\mathrm{Ker}}}
\def\Im{\mathop{\mathrm{Im}}}
\def\Coker{\mathop{\mathrm{Coker}}}
\def\Coim{\mathop{\mathrm{Coim}}}

% Boxtimes
%
\DeclareMathSymbol{\boxtimes}{\mathbin}{AMSa}{"02}

%
% Macros for moduli stacks/spaces
%
\def\QCohstack{\mathcal{QC}\!\mathit{oh}}
\def\Cohstack{\mathcal{C}\!\mathit{oh}}
\def\Spacesstack{\mathcal{S}\!\mathit{paces}}
\def\Quotfunctor{\mathrm{Quot}}
\def\Hilbfunctor{\mathrm{Hilb}}
\def\Curvesstack{\mathcal{C}\!\mathit{urves}}
\def\Polarizedstack{\mathcal{P}\!\mathit{olarized}}
\def\Complexesstack{\mathcal{C}\!\mathit{omplexes}}
% \Pic is the operator that assigns to X its picard group, usage \Pic(X)
% \Picardstack_{X/B} denotes the Picard stack of X over B
% \Picardfunctor_{X/B} denotes the Picard functor of X over B
\def\Pic{\mathop{\mathrm{Pic}}\nolimits}
\def\Picardstack{\mathcal{P}\!\mathit{ic}}
\def\Picardfunctor{\mathrm{Pic}}
\def\Deformationcategory{\mathcal{D}\!\mathit{ef}}
\newenvironment{reference}{}{}
\newenvironment{slogan}{}{}
\newenvironment{history}{}{}
\usepackage{cleveref}
\crefname{originallemma}{original lemma}{original lemmas}
\Crefname{originallemma}{Original lemma}{Original lemmas}
\makeatletter
\newlabel{lemma-relatively-pseudo-coherent}{{09VC}{source}{}{sourcepointer-1}{}}
\newlabel{lemma-relatively-pseudo-coherent@cref}{{[originallemma][0][]09VC}{[1][0][]source}{}{}{}}
\makeatother
\crefname{originallemma}{original lemma}{original lemmas}
\Crefname{originallemma}{Original lemma}{Original lemmas}
\makeatletter
\newlabel{modules-lemma-i-star-equivalence}{{08KS}{source}{}{sourcepointer-2}{}}
\newlabel{modules-lemma-i-star-equivalence@cref}{{[originallemma][0][]08KS}{[1][0][]source}{}{}{}}
\makeatother
\crefname{originallemma}{original lemma}{original lemmas}
\Crefname{originallemma}{Original lemma}{Original lemmas}
\makeatletter
\newlabel{morphisms-lemma-i-star-equivalence}{{01QY}{source}{}{sourcepointer-3}{}}
\newlabel{morphisms-lemma-i-star-equivalence@cref}{{[originallemma][0][]01QY}{[1][0][]source}{}{}{}}
\makeatother
\crefname{originallemma}{original lemma}{original lemmas}
\Crefname{originallemma}{Original lemma}{Original lemmas}
\makeatletter
\newlabel{perfect-lemma-pseudo-coherent}{{08E5}{source}{}{sourcepointer-4}{}}
\newlabel{perfect-lemma-pseudo-coherent@cref}{{[originallemma][0][]08E5}{[1][0][]source}{}{}{}}
\makeatother
\title{Relative pseudo-coherence forces quasi-coherent cohomology above the amplitude bound}
\author{The Stacks Project Authors}
\date{stacks-project:more-morphisms.tex:local-snapshot-20261011; GNU FDL 1.2 or later}
\begin{document}
\maketitle
\noindent\textit{Editorial note.} Statement span is lines 16703-16710 of sources/stacks-more-morphisms/pinned-original-source.tex; proof span is lines 16712-16732, the complete author proof (19 body lines).\par\medskip
\noindent\textit{Editorial note.} Navigation stubs: 4 cited-but-not-included labels are supplied as non-load-bearing pointers with their pinned tags where the snapshot tag index has them.\par\medskip
\noindent\textit{Editorial note.} extra\_packages is empty on purpose: the retained preamble already loads xy with options and a second declaration clashes in the assembler wrapper.\par\medskip
\noindent\textit{Editorial note.} Text is retained verbatim, including the author's own typos and placeholders.\par\medskip
\noindent\textit{Editorial note.} This lemma guarantees the statement of relative pseudo-coherence actually lands in the quasi-coherent derived category; complements L16070 and L16074 on base-change behaviour.\par\medskip
\noindent\textit{Allowed standard prerequisites (editorial).}\begin{itemize}
\item Explicit prerequisite boundary. The affine characterization of relative pseudo-coherence (same chapter, lemma-relatively-pseudo-coherence), the absolute pseudo-coherent cohomology criterion (Derived Categories of Schemes, lemma-pseudo-coherent), and the closed-immersion direct image equivalence of modules (Morphisms, lemma-i-star-equivalence; Modules, lemma-i-star-equivalence).
\end{itemize}\medskip
\ifdefined\hypertarget\hypertarget{sourcepointer-1}{}\fi
\noindent\textit{Source reference.} Original lemma 09VC: sources/stacks-more-morphisms/pinned-original-source.tex, source lines 16610--16622. Cited by the retained statement or proof but not part of the collected spans; supplied only as a navigation pointer. The argument does not depend on its content.\par\medskip
\ifdefined\hypertarget\hypertarget{sourcepointer-2}{}\fi
\noindent\textit{Source reference.} Original lemma 08KS: sources/stacks-more-morphisms/pinned-original-source.tex, source lines 16732--16732. Cited by the retained statement or proof but not part of the collected spans; supplied only as a navigation pointer. The argument does not depend on its content.\par\medskip
\ifdefined\hypertarget\hypertarget{sourcepointer-3}{}\fi
\noindent\textit{Source reference.} Original lemma 01QY: sources/stacks-more-morphisms/pinned-original-source.tex, source lines 16732--16732. Cited by the retained statement or proof but not part of the collected spans; supplied only as a navigation pointer. The argument does not depend on its content.\par\medskip
\ifdefined\hypertarget\hypertarget{sourcepointer-4}{}\fi
\noindent\textit{Source reference.} Original lemma 08E5: sources/stacks-more-morphisms/pinned-original-source.tex, source lines 16732--16732. Cited by the retained statement or proof but not part of the collected spans; supplied only as a navigation pointer. The argument does not depend on its content.\par\medskip
\begin{lemma}
\label{lemma-relative-pseudo-coherence}
Let $f : X \to S$ be a morphism of schemes which is locally of finite type.
If $E$ in $D(\mathcal{O}_X)$ is $m$-pseudo-coherent relative to $S$,
then $H^i(E)$ is a quasi-coherent $\mathcal{O}_X$-module for $i > m$.
If $E$ is pseudo-coherent relative to $S$, then $E$ is an object of
$D_\QCoh(\mathcal{O}_X)$.
\end{lemma}

\begin{proof}
Choose an affine open covering $S = \bigcup V_i$ and
for each $i$ an affine open covering $f^{-1}(V_i) = \bigcup U_{ij}$
such that the equivalent conditions of
Lemma \ref{lemma-relatively-pseudo-coherent}
are satisfied for each of the pairs $(U_{ij} \to V_i, E|_{U_{ij}})$.
Since being quasi-coherent is local on $X$, we may assume
that there exists an closed immersion $i : X \to \mathbf{A}^n_S$
such that $Ri_*E$ is $m$-pseudo-coherent on $\mathbf{A}^n_S$.
By Derived Categories of Schemes, Lemma \ref{perfect-lemma-pseudo-coherent}
this means that $H^q(Ri_*E)$ is quasi-coherent for $q > m$.
Since $i_*$ is an exact functor, we have $i_*H^q(E) = H^q(Ri_*E)$
is quasi-coherent on $\mathbf{A}^n_S$.
By Morphisms, Lemma \ref{morphisms-lemma-i-star-equivalence}
this implies that $H^q(E)$ is quasi-coherent as desired
(strictly speaking it implies there exists some quasi-coherent
$\mathcal{O}_X$-module $\mathcal{F}$ such that
$i_*\mathcal{F} = i_*H^q(E)$ and then
Modules, Lemma \ref{modules-lemma-i-star-equivalence}
tells us that $\mathcal{F} \cong H^q(E)$ hence the result).
\end{proof}

\end{document}
